% Einheit 6 von 8 – Situationsbäume (bank/zufallsexperimente-und-pfadregeln/e6.jsonl, zusammenbau v0.1)
%% TODO Zeitmarke Einheit 6: Marken-Zeile nicht lesbar
%% TODO Zweigzeile Teil 1: was hier gelernt wird (Satz in Schülersprache aus dem Einheitstitel) – Einheit 6
\einheitenkopf[e6]{Einheit 6 von 8 · Situationsbäume}
\zweigzeile{Abitur GK}
\setcounter{aufgabe}{24}

%% TODO Titel: Ich-kann-Satz fehlt in der Bank (Platzhalter: Kettenname) – Nr. 25
%% TODO Anweisung: ein Satz über der ersten Teilaufgabe fehlt in der Bank – Nr. 25
\begin{aufgabe}{Situationsbäume}
%% TODO \rechenplatz{n}: Schrittzahl des Grundfalls fehlt in der Bank (nur bei fester Schreibform, 2.3 b)
\begin{teile}
\teil „Unter den Radfahrern tragen 60\,\% einen Helm.“ Welcher Anteil?\\ \kreuz{Anteil aller (Pfad oder Rand)}\\ \kreuz{Anteil in einer Teilgruppe (Ast)}
\teil 30\,\% der Haushalte eines Orts haben einen Hund (H), die übrigen nicht (N). Unter den Hundehaltern haben 40\,\% auch eine Katze (K), unter den übrigen 20\,\%; O heißt ohne Katze. Beschrifte den Baum.

\baumzwei{H/,N/}{K/,O/}{K/,O/}
\teil Ergänze die fehlenden Äste. $P(A \cap X)$?

\baumzwei{A/0{,}6,B/}{X/0{,}25,Y/}{X/,Y/0{,}9}
\teil 35\,\% der Gäste eines Restaurants sind Kinder, unter ihnen bestellen 70\,\% Pommes. 26\,\% aller Gäste sind Erwachsene, die Pommes bestellen. Anteil der Pommes-Besteller unter den Erwachsenen? \leerfeld
\teil Maschine A fertigt 60\,\% der Teile mit 3\,\% Ausschuss, Maschine B den Rest mit 5\,\% Ausschuss. Weise nach: 3,8\,\% aller Teile sind Ausschuss.
\teil 70\,\% der Vereinsmitglieder sind Erwachsene, die übrigen Jugendliche. Unter den Erwachsenen kommen 40\,\% mit dem Rad zum Training, unter den Jugendlichen 80\,\%. Kommen mehr Erwachsene oder mehr Jugendliche mit dem Rad?
\teil Befragung zum Schwarzfahren: Bei 1 bis 4 auf dem Würfel beantwortet man F1 „Bist du schon schwarzgefahren?“, bei 5 oder 6 F2 „Bist du noch nie schwarzgefahren?“, jeweils ehrlich mit Ja (J) oder Nein (N). p ist der Anteil der Schwarzfahrer. Beschrifte den Baum.

\baumzwei{F1/,F2/}{J/,N/}{J/,N/}
\teil Vier Kugeln werden unabhängig voneinander je mit $\frac{1}{2}$ rot gefärbt, sonst weiß, und in einen Beutel gelegt. Zwei der vier Kugeln werden zufällig entnommen. P(beide rot)? \hfill (Abitur 2023 LK) \\ P = \leerfeld
\end{teile}
\end{aufgabe}

%% TODO Titel: Ich-kann-Satz fehlt in der Bank (Platzhalter: Kettenname) – Nr. 26
%% TODO Anweisung: ein Satz über der ersten Teilaufgabe fehlt in der Bank – Nr. 26
\begin{aufgabe}{Situationsbäume – Fehler finden, Begründen, Anwendung, Darstellungswechsel}
\begin{teile}
\teil 40\,\% der Befragten sind Frauen; 12\,\% aller Befragten sind Frauen mit Brille. Ben schreibt an den Ast „Brille unter den Frauen“ den Wert $0{,}12$. Finde den Fehler.
\teil Begründe: Die Wahrscheinlichkeit eines Merkmals der zweiten Stufe ist die Summe aller Pfade zu diesem Merkmal.
\teil Ein Onlineshop verschickt 25\,\% der Bestellungen per Express, davon kommen 4\,\% verspätet an; bei Standardversand 9\,\%. Wie viele von 2\,000 Bestellungen einer Woche kommen verspätet an? \leerfeld
\teil 45\,\% der Schüler sind in der Oberstufe (O), die übrigen nicht (U). In der Oberstufe haben 60\,\% ein Tablet (T), sonst 30\,\%; K heißt kein Tablet. Zeichne den Baum.

\baumzwei{O/,U/}{T/,K/}{T/,K/}
\end{teile}
\end{aufgabe}
